尽管 \method 展现了性能和成本收益，仍存在若干可供未来研究的局限性。

\paragraph{数据收集的可扩展性。}训练轮级路由器的主要瓶颈，是跨多个候选模型收集多样、长时程轨迹的成本。受计算预算约束，我们在 ScienceWorld 上评测最多 50 步的 episode，在 HLE 上评测最多 30 步的 episode。

\paragraph{缺乏在线适应。}目前，\method 运行于离线学习范式。尽管这显著降低训练成本，路由器却无法实时适应新颖或快速变化的环境。通过强化学习实现在线更新可能进一步提升稳健性，但与前沿模型持续在线交互的高昂成本，使其对当前学术研究仍具挑战。

\paragraph{提示词缓存与切换开销。}多轮路由固有的技术挑战，是在不同模型之间切换时可能丢失提示词缓存（例如 KV cache）。在许多基于 API 的部署中，频繁切换需要新模型重新处理整条轨迹前缀，从而可能增加延迟和逐轮成本。然而，我们的分析表明 \method 部分缓解了该问题：不同于在瞬时错误首次出现时就切换模型的反应式基线，\method 呈现更“从容”的切换模式，在更长的轮次序列中保持模型稳定性。如成本效率结果所反映，这种结构化的模型使用有助于平衡自适应模型选择与上下文缓存收益之间的权衡。
